\section{Tests und Validierung}
\label{sec:tests}

\subsection{Teststrategie}

Die Implementierung des Dispatchers wurde
durch \textbf{automatisierte Integrationstests} validiert.
Die Tests in \texttt{tests/integration\_test.py} testen das System
als Black-Box über die gRPC-Schnittstelle -- kein Mocking,
kein Unit-Test von Interna.

\noindent Begründung: Da gRPC die einzige externe Schnittstelle ist, geben
Integrationstests über diese Schnittstelle die stärkste
Sicherheit, dass das System korrekt funktioniert.
Unit-Tests würden interne Implementierungsdetails testen,
die sich ändern können ohne das Verhalten zu brechen.

\subsection{Testergebnisse Phase 2}

Alle Integrationstests wurden erfolgreich bestanden:

\begin{table}[H]
\centering
\begin{tabular}{llll}
\toprule
\textbf{Issue} & \textbf{Funktion} & \textbf{Tests} & \textbf{Ergebnis} \\
\midrule
\#13 & PostTask (Aufgabe einreichen)            & 6  & \textbf{6/6 PASS} \\
\#14 & Worker-Lookup (Namensdienst)             & 2  & \textbf{2/2 PASS} \\
\#15 & Dispatch (Task an Worker senden)         & 3  & \textbf{3/3 PASS} \\
\#16 & ReturnResult (Ergebnis empfangen)        & 9  & \textbf{9/9 PASS} \\
\#17 & GetResult (Ergebnis an Client)           & 8  & \textbf{8/8 PASS} \\
\#18 & Timeout-Handling + Retry-Logik           & 10 & \textbf{10/10 PASS} \\
\#19 & Strukturiertes Logging                   & 3  & \textbf{3/3 PASS} \\
\#20 & GET\_STATUS Monitoring                   & 4  & \textbf{4/4 PASS} \\
\#21 & Queue + Nebenläufigkeit                  & 4  & \textbf{4/4 PASS} \\
\midrule
     & \textbf{Gesamt}                          & \textbf{49} & \textbf{49/49 PASS} \\
\bottomrule
\end{tabular}
\caption{Integrationstestergebnisse Phase 2 (49/49 Tests bestanden)}
\end{table}

\subsection{Getestete Szenarien}

\subsubsection*{PostTask}
\begin{itemize}
  \item Task mit gültigem Typ wird akzeptiert, \texttt{task\_id} wird zurückgegeben.
  \item Task mit leerem \texttt{task\_type}: \texttt{INVALID\_ARGUMENT}.
  \item Task mit zu langem \texttt{task\_type} ($>$ 32 Zeichen): abgelehnt.
  \item Task mit zu großem Payload ($>$ 1024 Zeichen): abgelehnt.
  \item Mehrere Tasks werden unabhängig voneinander akzeptiert.
  \item Zurückgegebene \texttt{task\_id} ist numerisch und eindeutig.
\end{itemize}

\subsubsection*{ReturnResult}
\begin{itemize}
  \item COMPLETED: Zustand wechselt zu COMPLETED, Ergebnis wird gespeichert.
  \item FAILED: Zustand wechselt zu FAILED, Fehlermeldung gespeichert.
  \item Unbekannte \texttt{task\_id}: \texttt{Ack(success=False)}.
  \item Verspätetes Ergebnis nach FAILED: wird ignoriert (Idempotenz).
  \item \texttt{duration\_ms} wird korrekt berechnet.
  \item Timeout-Timer wird nach ReturnResult abgebrochen.
\end{itemize}

\subsubsection*{GetResult}
\begin{itemize}
  \item QUEUED-Task: Ergebnis leer, Status QUEUED.
  \item DISPATCHED-Task: Ergebnis leer, Status DISPATCHED.
  \item COMPLETED-Task: Ergebnis vorhanden, Status COMPLETED.
  \item FAILED-Task: Fehlermeldung, Status FAILED.
  \item Unbekannte \texttt{task\_id}: \texttt{NOT\_FOUND}.
  \item Polling-Verhalten: Status ändert sich korrekt im Zeitverlauf.
\end{itemize}

\subsubsection*{Timeout + Retry}
\begin{itemize}
  \item Task erhält nach Timeout den Zustand TIMEOUT.
  \item Retry: Task wechselt zu RETRYING und wird neu eingereiht.
  \item \texttt{retry\_count} wird bei jedem Timeout inkrementiert.
  \item Nach Erreichen von \texttt{MAX\_RETRIES}: Task wird FAILED.
  \item Timer wird bei rechtzeitigem ReturnResult abgebrochen (kein Spurious Timeout).
  \item Verspätetes Feuern auf terminalen Task: wird ignoriert.
\end{itemize}

\subsubsection*{Nebenläufigkeit}
\begin{itemize}
  \item 10 Tasks gleichzeitig eingereiht: alle werden korrekt verarbeitet.
  \item Keine Race Conditions bei gleichzeitigem Zustandswechsel.
  \item TaskStore und TaskQueue sind thread-sicher.
\end{itemize}

\subsection{Bekannte Testbesonderheiten}

\textbf{Encoding-Problem beim Ausführen auf Windows}:
Das Testskript gibt strukturierte Log-Ausgaben mit Sonderzeichen aus.
Unter Windows (PowerShell, cp1252-Encoding) führte dies zu
\texttt{UnicodeEncodeError}.
Lösung: \texttt{\$env:PYTHONIOENCODING = "utf-8"} vor dem Testaufruf setzen.

\noindent \textbf{Doppelte Log-Ausgaben}:
Vor der Einführung von \texttt{logger.propagate = False} wurde jeder
Log-Eintrag mehrfach ausgegeben (Python-Logger-Hierarchie).
Dieser Fehler wurde in \texttt{src/common/logger.py} behoben.
Die 49/49 Tests bestehen weiterhin nach dem Fix.

\subsection{Phase 3 -- Docker-Tests}

Die Docker-Konfiguration wurde durch manuelle
Überprüfung aller Dockerfile- und docker-compose.yml-Konfigurationen
validiert.
Folgende Aspekte wurden geprüft:

\begin{itemize}
  \item Layer-Cache-Strategie in allen Dockerfiles korrekt.
  \item CRLF/LF-Konvertierung für \texttt{entrypoint.sh} implementiert.
  \item Alle Umgebungsvariablen korrekt gemappt
        (insbesondere \texttt{WORKER\_TYPE} $\rightarrow$ \texttt{TASK\_TYPES}).
  \item Health Checks syntaktisch korrekt und funktional.
  \item \texttt{depends\_on} mit \texttt{condition: service\_healthy} korrekt.
  \item Kein Task/Ergebnis-Volume vorhanden (§8-Konformität).
  \item Worker-Ports nur intern (kein externes Mapping).
\end{itemize}

Ein vollständiger End-to-End-Test des containerisierten Systems steht aus, da die Namensdienst-
und Worker-Implementierungen von anderen Teammitgliedern noch
nicht abgeschlossen sind.
